\section{组织：AI 全站工程师、远程协作和高人才密度}

前面讨论产品和平台，本章看 ONE2X 的组织方式。王冠把 ONE2X 定位成 AI 时代的产品工作室，而不是传统公司。工作室强调探索、兴趣、质量和超级个体。团队里每个人要能独当一面，被称为 AI 全站工程师。

\subsection{为什么小团队能做更多}

本节从平台理论回到组织执行。生成系统不是只靠宏大叙事，需要一个能快速验证、快速产出、快速学习的小团队来把方法落地。

AI 时代很多生产能力以 token 形式存在，等于团队可以调用外部劳动力。高人才密度小团队如果有强自驱力和清晰方向，就能做出过去需要更大团队才能完成的产出。ONE2X 的招聘重点包括自驱力、主动性、公开构建、开源项目、认知分享，以及个人 reward function 是否和团队方向有投影。

\begin{knowledgebox}{组织关键词}
\begin{tabularx}{\textwidth}{p{0.22\textwidth}p{0.34\textwidth}X}
\toprule
关键词 & 含义 & 组织作用 \\
\midrule
AI 全站工程师 & 能跨产品、工程、模型、内容和业务理解问题的人 & 适合探索期的小团队。 \\
超级个体 & 高自驱、能独立产出、能用 AI 放大能力的人 & 降低管理和协作成本。 \\
Reward function & 个体真正被什么激励 & 招聘时要和团队目标同向。 \\
Remote & 远程办公 & 扩大人才池，但需要解决孤独和信任。 \\
\bottomrule
\end{tabularx}
\end{knowledgebox}

\subsection{远程办公的问题与解法}

远程办公带来孤独感和信任问题，因为文字留言缺少语气、肢体和上下文。王冠没有因此放弃 remote，而是把它当成一种组织约束：用文档、异步沟通、强自驱和明确目标来抵消问题。线下办公也有线下办公的问题，远程不是更轻松，而是用另一套方式换取更大人才池和更高自主性。

\subsection{本章小结}

ONE2X 的组织形态服务于生成系统探索：小团队、高人才密度、强自驱、远程协作、AI 放大生产力。它不是普适答案，但和“产品工作室”定位一致。

